\sectionkicker{Theme synthesis}
\subsection*{横断比較 — Multimodalからinteraction/action spaceへ}
\addcontentsline{toc}{subsection}{横断比較 — Multimodalからinteraction/action spaceへ}
video generation、world model、roboticsは同じmultimodalという語で括れてもavailabilityと目的が異なる。2025年後半は生成modalitiesの増加だけでなく、環境理解・予測・physical actionへ接続する研究と製品面が並行した。
\medskip
\begin{center}
\small
\renewcommand{\arraystretch}{1.16}
\begin{tabularx}{\linewidth}{@{}p{0.20\linewidth}X@{}}
\toprule
観察軸 & 一次資料・論文から読めること \\
\midrule
interactive world / robotics & DeepMindのGenie 3、Gemini Robotics 1.5、SIMA 2などはinteractive world、robotics、3D agentを別eventとして展開した。research prototype、partner access、API availabilityを一つのrelease状態へ統合しない。 \cite{src-dd5d9b8662f3} \\
audio-video generation & Sora 2はvideo generationにaudioを含むrelease identityを持つ。生成modalitiesの拡大はroboticsやworld-model研究と関連づけて読めるが、同一product categoryではない。 \cite{src-661a3476d3d5,src-8530f5b0e216} \\
model/product family & MiniMaxのM2、Hailuo 2.3、M2.1は同社内でもmodel/product surfaceが異なる。client-rendered indexの境界を保ち、item-level release identityを使って整理する。 \cite{src-d8e13386974c} \\
availabilityの読み分け & この章の各eventはresearch announcement、limited access、managed API、product rolloutなど提供状態が異なる。技術方向が近いことと、同じ条件で利用可能であることを混同しない。 \cite{src-dd5d9b8662f3,src-661a3476d3d5,src-8530f5b0e216,src-d8e13386974c} \\
\bottomrule
\end{tabularx}
\renewcommand{\arraystretch}{1.0}
\end{center}
\normalsize
\begin{claimboundary}[この比較の境界]
この表は本号で既に検証・参照している一次資料と論文を横断比較の形へ再配置したものである。vendor / project / author が報告した性能値や優位性は独立再現済みの一般則へ変換せず、本文とSource Notesの帰属境界をそのまま維持する。
\end{claimboundary}
